\scanpage{26}
On trouve ainsi une bijection canonique
\[
T\simeq\{\mathrm{tr}\,\bar a\}\amalg a\times\varphi,
\]
évidemment compatible avec l'action de $G_s\simeq\unsure{\mathfrak S|\mathscr S}_a\times\mathfrak S_\varphi\simeq\mathbb Z/2\mathbb Z\times\mathbb Z/2\mathbb Z$, qui opère donc \uline{fidèlement} sur $T$. On retrouve que $G_s$ opère sur $T$ par permutations paires, et \deleted{bref} de plus on conclut que l'image de $G$ dans $\mathrm{Alt}_T$ est d'ordre divisible par $4$. Étant déjà \inserted{d'ordre} divisible par $30$, il est divisible par $60$ Or $\mathrm{Alt}_T\deleted{\unsure{\{}}$ est d'ordre $60$, \deleted{on trouve} et \deleted{Aut} $G$ d'ordre $120$, on trouve

\uline{Th} L'hom canonique
\[
G=\mathrm{Aut}\underset{\substack{|\\\text{icos. comb}}}{\Pi}
 \longrightarrow\mathfrak S_{\underbrace{T(\Pi)}_{\substack{\text{ens des}\\\text{trièdres distingués}}}}
\]
a comme image $\mathrm{Alt}_T$, et comme noyau le ss-groupe \inserted{$\mathfrak z$} d'ordre $2$ engendré par l'antipodisme $\underline a$.

Comme \deleted{\illegible{1}} $\Pi$ \uline{est orientable} et $\underline a$ renverse \deleted{Corollaire} l'orientation, on trouve
\[
G\simeq \mathfrak z\times G^+
\]
donc
\[
G^+\xrightarrow{\sim}G/\mathfrak z
\]
donc
% archive sheet number: 25
